\appendix
\section{Detailed Dynamic Pseudo-Regret Analysis}
\label{app:proof}

This appendix proves the bound in Section~\ref{sec:theory}. The result is conditional on the time-varying linear delta model and the window-coverage condition stated there. These assumptions are necessary: a changing latent profile alone does not establish a non-stationary bandit regret bound.

\subsection{Filtration and confidence event}

Let \(\mathcal{F}_{t-1}\) be the sigma-field generated by all observations and selected tasks before bandit round \(t\). For the selected task, abbreviate \(\boldsymbol{\phi}_{t}=\boldsymbol{\phi}_{t,\tau_t}\). Assumption~1 can be written as
\begin{equation}
    \boldsymbol{\delta}_{t}=\mathbf{W}_t^{*\top}\boldsymbol{\phi}_{t}+\boldsymbol{\varepsilon}_{t},
    \qquad
    \mathbb{E}[\boldsymbol{\varepsilon}_{t}\mid\mathcal{F}_{t-1},\tau_t]=\mathbf{0}.
    \label{eq:appendix-linear-model}
\end{equation}
For \(\mathbf{u}\in\mathbb{R}^{K}\) with \(\lVert\mathbf{u}\rVert_2\leq1\), the scalar noise \(\mathbf{u}^{\top}\boldsymbol{\varepsilon}_t\) is conditionally \(\sigma\)-sub-Gaussian.

Recall the ridge matrices formed from completed observations in \(\mathcal{W}_t\):
\begin{equation}
    \mathbf{A}_{t}=\lambda\mathbf{I}_{d_f}+\sum_{i\in\mathcal{W}_t}\boldsymbol{\phi}_{i}\boldsymbol{\phi}_{i}^{\top},
    \qquad
    \mathbf{B}_{t}=\sum_{i\in\mathcal{W}_t}\boldsymbol{\phi}_{i}\boldsymbol{\delta}_{i}^{\top},
    \qquad
    \widehat{\mathbf{W}}_t=\mathbf{A}_t^{-1}\mathbf{B}_t.
    \label{eq:appendix-ridge}
\end{equation}
For a failure probability \(\rho\in(0,1)\), define
\begin{equation}
    \beta_t=\sigma\sqrt{2K\log\!\left(
    \frac{K\det(\mathbf{A}_t)^{1/2}}{\rho\lambda^{d_f/2}}
    \right)}+\sqrt{\lambda}S.
    \label{eq:appendix-beta}
\end{equation}
The factor \(K\) follows from applying a scalar self-normalized concentration inequality to each output coordinate and taking a union bound. It can be replaced by a sharper vector-valued confidence radius when additional noise structure is available.

\subsection{Estimation error and drift bias}

For a candidate \(\tau\), let
\[
    w_t(\tau)=\sqrt{\boldsymbol{\phi}_{t,\tau}^{\top}\mathbf{A}_t^{-1}\boldsymbol{\phi}_{t,\tau}}.
\]
The following lemma separates stochastic estimation error from bias due to stale samples.

\paragraph{Lemma 1 (windowed prediction error).}
Under Assumptions~1--4, with probability at least \(1-\rho\), simultaneously for every \(t\leq T\) and \(\tau\in\mathcal{T}\),
\begin{equation}
\left|
\mathbf{g}_t^{\top}(\widehat{\mathbf{W}}_t-\mathbf{W}_t^*)^{\top}
\boldsymbol{\phi}_{t,\tau}
\right|
\leq
\beta_t\lVert\mathbf{g}_t\rVert_2w_t(\tau)+D_t,
\label{eq:appendix-prediction-error}
\end{equation}
where, for a constant \(c_{\lambda,\kappa}>0\) depending only on the regularization and coverage constants,
\begin{equation}
    D_t\leq c_{\lambda,\kappa}
    \sum_{j=\max\{1,t-\ell\}}^{t-1}
    \lVert\mathbf{W}_{j+1}^{*}-\mathbf{W}_{j}^{*}\rVert_F.
    \label{eq:appendix-drift-bias}
\end{equation}

\paragraph{Proof.}
Substituting~\eqref{eq:appendix-linear-model} into~\eqref{eq:appendix-ridge} and adding and subtracting \(\mathbf{W}_t^*\) gives
\begin{align}
\widehat{\mathbf{W}}_t-\mathbf{W}_t^*
={}&\mathbf{A}_t^{-1}\sum_{i\in\mathcal{W}_t}
\boldsymbol{\phi}_i\boldsymbol{\varepsilon}_i^{\top}
-\lambda\mathbf{A}_t^{-1}\mathbf{W}_t^* \nonumber\\
&+\mathbf{A}_t^{-1}\sum_{i\in\mathcal{W}_t}
\boldsymbol{\phi}_i\boldsymbol{\phi}_i^{\top}
\left(\mathbf{W}_i^*-\mathbf{W}_t^*\right).
\label{eq:appendix-error-decomposition}
\end{align}
The first two terms are the usual stochastic and regularization terms of ridge regression. Applying the self-normalized inequality coordinatewise, followed by Cauchy--Schwarz over the \(K\) output coordinates, bounds their scalarized prediction error by
\[
\beta_t\lVert\mathbf{g}_t\rVert_2
\lVert\boldsymbol{\phi}_{t,\tau}\rVert_{\mathbf{A}_t^{-1}}.
\]
For the final term, Assumption~4 bounds the inverse window Gram matrix, while \(\lVert\boldsymbol{\phi}_i\rVert_2\leq1\). Thus, its scalarized prediction error is at most a constant \(c_{\lambda,\kappa}\) times
\[
\sum_{i\in\mathcal{W}_t}\lVert\mathbf{W}_i^*-\mathbf{W}_t^*\rVert_F.
\]
By telescoping \(\mathbf{W}_i^*-\mathbf{W}_t^*\) along the parameter path, this is bounded by the right-hand side of~\eqref{eq:appendix-drift-bias}. Combining the three bounds proves the lemma. \hfill\(\square\)

Summing~\eqref{eq:appendix-drift-bias} over time counts each increment \(\lVert\mathbf{W}_{j+1}^*-\mathbf{W}_j^*\rVert_F\) in at most \(\ell\) windows. Hence,
\begin{equation}
    \sum_{t=1}^{T}D_t\leq c_{\lambda,\kappa}\ell P_T.
    \label{eq:appendix-drift-sum}
\end{equation}

\subsection{From confidence to dynamic pseudo-regret}

Define the confidence width
\[
    U_t(\tau)=\beta_t\lVert\mathbf{g}_t\rVert_2w_t(\tau).
\]
On the event in Lemma~1, the expected reward satisfies, for every candidate \(\tau\),
\begin{equation}
    q_t(\tau)\leq
    \mathbf{g}_t^{\top}\widehat{\mathbf{W}}_t^{\top}\boldsymbol{\phi}_{t,\tau}
    +U_t(\tau)+D_t.
    \label{eq:appendix-upper-confidence}
\end{equation}
The selected task maximizes the first two terms on the right-hand side. Applying~\eqref{eq:appendix-upper-confidence} to \(\tau_t^*\), then applying the corresponding lower bound from Lemma~1 to \(\tau_t\), yields
\begin{align}
q_t(\tau_t^*)-q_t(\tau_t)
&\leq
\left[\mathbf{g}_t^{\top}\widehat{\mathbf{W}}_t^{\top}\boldsymbol{\phi}_{t,\tau_t}
+U_t(\tau_t)+D_t\right]-q_t(\tau_t) \nonumber\\
&\leq 2U_t(\tau_t)+2D_t.
\label{eq:appendix-instant-regret}
\end{align}
This is approximate optimism: the stale-sample term \(D_t\) prevents exact UCB optimism, but is explicitly controlled by the parameter-path variation.

Because \(\mathbf{g}_t\) is a softmax vector, \(\lVert\mathbf{g}_t\rVert_2\leq1\). Once the window is full, Assumption~4 gives \(\mathbf{A}_t\succeq\kappa\ell\mathbf{I}_{d_f}\). The warmup initialization supplies a full window before the first bandit decision. Consequently,
\begin{equation}
    \sum_{t=1}^{T}U_t(\tau_t)
    \leq\frac{\beta T}{\sqrt{\kappa\ell}},
    \qquad
    \beta=\max_{t\leq T}\beta_t.
    \label{eq:appendix-width-sum}
\end{equation}
Summing~\eqref{eq:appendix-instant-regret}, and using~\eqref{eq:appendix-drift-sum} and~\eqref{eq:appendix-width-sum}, proves
\[
R_T\leq
O\!\left(\frac{\beta T}{\sqrt{\kappa\ell}}\right)
+O(\ell P_T).
\]
Balancing the two terms gives the window choice and regret rate in Section~\ref{sec:theory}. \hfill\(\square\)
